% ECONOMICS_PROBLEM: {"id": "D72-527", "title": "Information that creates electoral accountability", "jel_code": "D72", "source_id": "OP-0387"}
% !TeX program = xelatex
\documentclass[11pt,a4paper]{article}
\usepackage[margin=23mm]{geometry}
\usepackage{fontspec}
\setmainfont{Latin Modern Roman}
\usepackage{amsmath,amssymb,mathtools}
\usepackage{xcolor,enumitem,booktabs,longtable,array,xurl,hyperref}
\definecolor{muted}{HTML}{53636C}
\hypersetup{colorlinks=true,urlcolor=blue,linkcolor=blue}
\setlength{\parindent}{0pt}
\setlength{\parskip}{5pt}
\newcommand{\R}{\mathbb R}
\newcommand{\E}{\mathbb E}
\newcommand{\N}{\mathbb N}
\newcommand{\ind}{\mathbf 1}
\newcommand{\Var}{\operatorname{Var}}
\newcommand{\Cov}{\operatorname{Cov}}
\newcommand{\supp}{\operatorname{supp}}
\DeclareMathOperator*{\argmin}{arg\,min}
\DeclareMathOperator*{\argmax}{arg\,max}
\newcommand{\entrymeta}[3]{{\sffamily\small\color{muted}\textbf{\texttt{#1}}\quad|\quad #2\par #3\par}}
\newcommand{\Problemref}[1]{\texttt{#1}}
\newcommand{\JELref}[2]{\texttt{#2} (JEL \texttt{#1})}
\begin{document}
\section*{Information that creates electoral accountability}
\textbf{Problem ID:} \texttt{D72-527}\quad \textbf{JEL:} \texttt{D72}\par
% BEGIN ECONOMICS STATEMENT
\label{problem:OP-0387}
{\sffamily\small\color{muted}Original subject group: Political economy and institutions\par}
\label{definitions:problem:30007669}
\textbf{Audit classification:} Background; proposed extension. All 29 authors, year and DOI checked. PMC was blocked by a challenge; accessible article abstract supplies null average behavior findings, not a certified long-run services conjecture.
\subsection*{Definitions and precise research target}
\textbf{Complete editorial problem.} Fix an electoral setting and a verified measure \(q_j\) of incumbent \(j\)'s public-service performance. An information intervention \(z\) delivers a declared message through a specified channel before an election. Let \(B_i(z)\) be voter \(i\)'s belief about performance, \(V_i(z)\) the vote choice, and \(G_{jT}(z)\) subsequent service provision. Define belief, vote, and service effects separately, including \(\tau_G=E[G_{jT}(1)-G_{jT}(0)]\). Accountability additionally requires evaluating how support responds to performance, rather than an undirected change in incumbent vote share. Orthogonal message, source-credibility, coordination, or alternative-candidate variations can distinguish information failures from distrust, identity preferences, clientelistic incentives, and inability to coordinate. State assumptions permitting those mechanism interpretations. Account for message spillovers, selective exposure, turnout, and politicians' responses to anticipated disclosure. A null average vote effect can coexist with corrected beliefs or subgroup effects; it does not show that information never matters. The proposed problem is to identify conditions yielding durable accountability and service improvements. Existing coordinated trials provide substantial answers for their tested designs; an exact primary open-question paragraph for the combined catalogue formulation was not verified.

\textbf{Definitions and admissible scope.} Define \(V_i\in\{0,1,\varnothing\}\) as incumbent support, alternative support, or nonparticipation; assigning turnout a separate code avoids conditioning on post-treatment voters. Verified incumbent performance \(q_j\) is measured before messages and on a common scale. For a binary high/low performance category, an accountability contrast is \(\kappa=[E(V^I(1)-V^I(0)\mid q=H)]-[E(V^I(1)-V^I(0)\mid q=L)]\), where \(V^I\) is the incumbent-vote indicator among all eligible voters. It is heterogeneity of information effects, not a causal effect of changing performance. Later service \(G_{jT}\) is measured at jurisdiction level and may depend on election outcomes and politicians' anticipation. Belief, vote, turnout, and service estimands are distinct. Take a fixed eligible-voter cohort and jurisdiction assignment law. Define \(V_i^I(z)=1\{V_i(z)=1\}\), so nonvoters contribute zero rather than disappearing. Belief effects use a fixed score, for example the mean posterior probability of high performance; vote and service effects use their own scales. Specify dates of information, election, and service follow-up \(T\), source accuracy and delivery, performance categories, and cluster exposure mappings. The source's existing average vote evidence is a constraint on relevant model interpretations, while the joint identified set of belief, turnout, accountability-heterogeneity, and later-service effects is the editorial target. Subgroup contrasts require common design support and no claim that observed performance was randomized.
\subsection*{Variants and assumption changes}
\begin{enumerate}
\item Typical message replication (source-based scope): replicate comparable nonpartisan performance messages with preregistered pooled effects. The coordinated trials already estimate average voter behavior effects; their null finding is evidence, not proof that every information design fails.
\item Coordination and credible alternatives (editorial): independently vary trusted source, community discussion, and alternative-candidate information and follow subsequent service provision. Cluster assignment and politician spillovers must be defined before mechanism claims.
\end{enumerate}
\subsection*{References and source audit}
\textbf{Support and access.} All 29 authors, year and DOI checked. PMC was blocked by a challenge; accessible article abstract supplies null average behavior findings, not a certified long-run services conjecture. \textbf{Locator:} Article abstract reproduced in PubMed; Science Advances 5(7), eaaw2612.
\begin{enumerate}\raggedright
\item\hypertarget{definitions:source:30007669:1}{} Thad Dunning; Guy Grossman; Macartan Humphreys; Susan D. Hyde; Craig McIntosh; Gareth Nellis; Claire L. Adida; Eric Arias; Clara Bicalho; Taylor C. Boas; Mark T. Buntaine; Simon Chauchard; Anirvan Chowdhury; Jessica Gottlieb; F. Daniel Hidalgo; Marcus Holmlund; Ryan Jablonski; Eric Kramon; Horacio Larreguy; Malte Lierl; John Marshall; Gwyneth McClendon; Marcus A. Melo; Daniel L. Nielson; Paula M. Pickering; Melina R. Platas; Pablo Querubín; Pia Raffler; Neelanjan Sircar. 2019. \emph{Voter information campaigns and political accountability: Cumulative findings from a preregistered meta-analysis of coordinated trials}. Article abstract and bibliographic metadata; Science Advances 5(7), eaaw2612.
\url{https://pubmed.ncbi.nlm.nih.gov/31281891/}.
DOI: \url{https://doi.org/10.1126/sciadv.aaw2612}.
\end{enumerate}

\subsection*{Common conventions: Source mathematical conventions}

These conventions apply to each entry unless explicitly replaced.

\textbf{Models and identification.} A primitive model \(m\) specifies all probability laws, feasible choices, information, equilibrium and measurement rules used by its target. The admissible class \(\mathcal M\) is fixed independently of outcomes. If \(P_O(m)\) is its observable probability law and \(\tau(m)\) the target, its identified set at observable law \(P\) is
\[
 \mathcal I_\tau(P)=\{\tau(m):m\in\mathcal M,\ P_O(m)=P\}.
\]
Point identification means that this set is a singleton. An empty set signals incompatibility of data and assumptions. Coordinate bounds need not describe the joint set. Bounds are sharp only if every retained target is attainable or arbitrarily approachable by compatible models. Finite bounds require explicit restrictions on unobserved quantities; unrestricted confounding, hidden wealth, extrapolation or arbitrary equilibrium selection can yield uninformative sets.

\textbf{Causal interventions.} A potential outcome \(Y_i(p)\) is the outcome under a complete policy regime \(p\), including timing, financing, information and market spillovers. The target population and units are fixed before treatment. With interference, use the complete assignment vector or a justified exposure mapping; individual no-interference notation alone cannot identify an economy-wide intervention. A difference of mean potential outcomes does not require the cross-world joint distribution. Conditioning on post-treatment migration, survival, enrollment or employment changes the estimand and needs separate justification.

\textbf{Statistical statements.} An estimator, confidence set, forecast or test specifies a sampling law, model class and loss/coverage criterion. Uniform \((1-\alpha)\)-coverage means \(\inf_{m\in\mathcal M}\Pr_m\{\tau(m)\in C_n\}\geq1-\alpha\), or an explicitly asymptotic version. The whole parameter space trivially has coverage, so informativeness requires a stated width, risk or power objective. A forecasting improvement is relative to a named benchmark, information set and evaluation population.

\textbf{Equilibrium and optimization.} An equilibrium comprises feasible plans, optimality, beliefs where needed, budget constraints and all market/resource clearing. The primitives or a stated benchmark define each correspondence \(\mathcal E(p;m)\). Nonemptiness is assumed or proved, never inferred just from compact policy sets. A selected-equilibrium target uses a declared selection rule; a robust target quantifies over all permitted equilibria. Use suprema or infima unless attainment is established. An \(\varepsilon\)-optimal policy is feasible and within \(\varepsilon\) of the finite optimum value. Report an empty feasible set separately.

\textbf{Welfare and accounting.} Normative utility, interpersonal normalization, social weights, numeraire, discounting and the use of public receipts are primitives. Behavioral utility need not coincide with the welfare criterion. Household welfare includes ownership income and rents once; adding profits again to compensating variation generally double-counts them. A monetary equivalent is defined by an explicit transfer and price convention, with monotonicity and range conditions. Average effects, mechanism contributions, transfers and welfare losses are different targets. Infinite sums require convergence and dynamic feasible plans require terminal or no-Ponzi conditions.

\textbf{Source versions.} A working-paper date, revision date and journal publication year are distinguished. A DOI for a journal version is not automatically the DOI of an earlier manuscript. Locators refer to the stated version and printed pages where specified. A metadata-only check does not verify a claim about a full-text conclusion. Every entry's source subsection narrows the provenance claim accordingly.
% END ECONOMICS STATEMENT
\end{document}
